\section{Discussion and Conclusion}
\label{sec:discussion}

The sequence developed here isolates a form of order that remains stable while saturated with predation. Equality survives there only in the thin accounting sense that every loss is offset by a corresponding gain, and once that accounting is disturbed the successor need not be cleaner, more egalitarian, or morally mixed in any interesting way. It can instead become a harder arrangement in which the distribution created during the transition is defended more effectively than it was created. The stages of that movement remain distinct: balanced reciprocity governs the unit benchmark, a permanent honest perturbation destroys stationarity once exposure becomes wealth dependent, delegated predation emerges when wealth and vulnerability diverge enough to create surplus from hiring, and organized protection arrives when concentrated wealth becomes worth defending and the prey base has already been thinned. Welfare follows the same sequence. Honesty does not become socially harmful; welfare falls because a local rupture in a corrupt equilibrium is followed by unequal adjustment when no broader institutional change accompanies it. Some agents reduce exposure by staying home, then by hiring thieves, and later by hiring guards, while others become more exposed as their outside options weaken with repeated loss. Concavity in wealth, theft effort costs, contracting costs, guarding costs, and punishment losses all move in the same direction.

Guards and police appear only after concentrated wealth has become worth protecting and labor can be shifted from predation to protection without altering who controls the terms of social order. The move to guarding therefore replaces one order with another; it does not mark the first appearance of order. Simulation is useful precisely because that replacement unfolds through state dependence that is hard to characterize in closed form once many heterogeneous agents interact: exact event time makes the first honest night visible, patronage and enforcement paths show how the payoff comparisons behave after symmetry has broken, and the phase diagram together with the robustness panels show where the mechanism survives and where it gives way to coexistence or collapse. Production is absent, so the analysis concerns the redistribution, dissipation, and protection of wealth rather than its endogenous creation, and the final regime remains protection of accumulated spoils rather than a full productive property-rights order. The case studied is also the one that gives the fable its force, namely a small honest rupture inside a society whose predatory balance relied on universal participation, while a broader coalition of honest agents could generate a different transition.

The paradox that made the story worth modeling remains: universal theft can look balanced, honesty can break that balance without redeeming the society, and the calmer order that follows can also be the more unequal order. Those claims now rest on explicit conditions. A balanced theft cycle can persist under symmetry, a permanent honest perturbation destroys stationarity once exposure becomes wealth dependent, patronage requires surplus over poor outside options, and guarding becomes privately optimal once concentrated wealth and prey depletion change the ranking of control technologies. The computational evidence completes the sequence by showing a precise first night asymmetry, a durable widening of the distribution, a contract regime grounded in explicit surplus, and a later move toward protection that stabilizes gains accumulated during the predatory transition. Calvino remains central because the story locates the whole movement with unusual economy.
